\documentclass[11pt]{article}
\usepackage[margin=1in]{geometry}
\usepackage[T1]{fontenc}
\usepackage{lmodern}
\usepackage{microtype}
\usepackage{amsmath}
\usepackage{booktabs}
\usepackage[hidelinks]{hyperref}
\usepackage[sort&compress,numbers]{natbib}

\newcommand{\DarkNESS}{\textsc{DarkNESS}}
\newcommand{\BpL}{B_{\perp}L}
\newcommand{\Th}{\Theta}
\newcommand{\gag}{g_{a\gamma\gamma}}
\newcommand{\todo}[1]{\textbf{[#1]}}

\title{A Skipper-CCD NanoSat search for geomagnetic conversion\\ of axion-like particles from dark matter decay}
\author{[Author name]}
\date{\today}

\begin{document}
\maketitle

\begin{abstract}
A long-lived dark matter parent decaying into two relativistic axion-like particles (ALPs) yields a narrow line at half the parent mass from the Milky Way halo and a redshifted continuum below it from the universe. An ALP crossing Earth's magnetic field converts to an X-ray photon with a probability set by the ALP--photon coupling and the transverse field integrated along the line of sight, so the converted brightness is proportional to $\Th=\gag^{2}f/\tau$ times a kernel that orbit and pointing determine. We derive this channel from the interaction term to the counts in a \DarkNESS{}-class skipper-CCD NanoSat, check each step against closed forms and the published Suzaku analysis, and evaluate the search with a Poisson likelihood profiled over celestial, particle and Earth-limb backgrounds. On the reference observing schedule the 90\% statistical limit on $\Th$ lies $4\times10^{5}$ to $2\times10^{6}$ above the ceiling allowed by the stellar bound on $\gag$ and the Planck 2018 bound on decaying dark matter; no allowed model is detectable and the gap does not close with exposure. The analysis establishes instead how such a search must be built. The signal is identified by its energy spectrum, not its orbital modulation: a time-only fit keeps 6\% of the information because the limb angle and cutoff rigidity track the conversion kernel, and a time--energy fit keeps 95\%. The halo line carries the sensitivity; the continuum alone is 37 times weaker. The same machinery reproduces the statistical floor of the Suzaku analysis within a factor of two of its published coupling limit.
\end{abstract}

\noindent\textbf{Keywords:} axion-like particles; dark matter; Skipper-CCD; NanoSatellite; diffuse X-rays; mission design; geomagnetic field

\section{Introduction}
\label{sec:introduction}

The published \DarkNESS{} observatory establishes a compact platform for low-energy X-ray measurements from low Earth orbit (LEO). Its reference payload combines four fully depleted Skipper-CCDs, 12~cm$^2$ of geometric sensor area, a direct-view aperture without X-ray optics, a $20^\circ$ field of view, and a nominal 1--10~keV response \citep{alpine2025darkness}. This work treats that payload and spacecraft as an observatory architecture rather than as a fixed observing program.

The objective considered is a search for a specific dark matter decay channel. A long-lived, non-relativistic parent particle $\chi$, comprising all or a stated fraction of the dark matter, decays into two relativistic ALPs. For a parent at rest each daughter carries half the parent mass; parent masses of 2--20~keV place the daughters in the \DarkNESS{} band. Galactic decays produce a direction-dependent line near that energy, while decays over cosmological distances produce a redshift-broadened continuum from the same process \citep{yamamoto2020alp,dror2021cosmic,cui2022cosmogenic}. The ALP is a relativistic messenger from dark matter decay, not the local cold dark matter.

The ALPs are not detected through an energy deposition of their own. Their coupling to electromagnetism permits conversion to photons in Earth's magnetic field through the inverse Primakoff effect. In the weak-mixing regime the conversion amplitude depends on the field component transverse to the trajectory and on the phase accumulated along it \citep{sikivie1983experimental,raffelt1988mixing}. The converted photon keeps the ALP's energy and direction. For a detector in LEO the conversion probability changes with position and line of sight, which supplies the defining observable: a response-folded X-ray component whose time dependence follows a kernel calculated from the orbit, attitude, field of view and geomagnetic field. A candidate signal must have the energy dependence of the decay model and the time dependence of geomagnetic conversion; an unexplained absolute excess is not evidence.

The Suzaku study of Yamamoto et al.\ is the observational benchmark \citep{yamamoto2020alp}. It calculated the conversion strength at 60-second cadence for four blank fields and fitted the spectra against it, found no dependence, and set a 99\% limit of $1.6\times10^{-9}\ \mathrm{erg\,s^{-1}\,cm^{-2}\,sr^{-1}}$ in 2--6~keV at $(\BpL)^2=10^{4}\ \mathrm{T^2\,m^2}$. It also named the central difficulty: conversion strength and charged-particle background vary with the same orbital variables.

This paper derives the channel independently of that analysis, from the interaction term to the counts in a \DarkNESS{}-class detector, and asks two questions of it: whether any model allowed by current constraints is within reach, and, given the answer, what the derivation establishes about how the measurement must be made. Section~\ref{sec:physics} derives the conversion and the source, Section~\ref{sec:detector} the counts, Section~\ref{sec:likelihood} the likelihood and the information that survives the backgrounds, and Section~\ref{sec:results} states the result.

\section{The channel}
\label{sec:physics}

\subsection{Conversion in the geomagnetic field}

From the interaction $-\tfrac14 g\,a F\tilde F = g\,a\,\mathbf E\cdot\mathbf B$ \citep{raffelt1988mixing}, linearised about a static field $\mathbf B_0$ for a plane wave of energy $\omega$ along $z$ and reduced to first order in $z$ for a slowly varying envelope, the ALP and the two transverse photon polarisations obey a Schr\"odinger-like system with mixing terms $\Delta_j = gB_j/2$ ($j=x,y$) and diagonal terms $\Delta_a=-m_a^2/2\omega$, $\Delta_{\rm pl}=-\omega_{\rm pl}^2/2\omega$ \citep{marsh2022fourier}. To first order in $g$ the amplitude from an ALP to polarisation $j$ over a path of length $L$ is $-\tfrac{ig}{2}\int_0^L B_j(s)\,e^{-iqs}ds$, and the two polarisations add incoherently:
\begin{equation}
P_{a\to\gamma}(E,\hat n,t)=\left(\frac{g}{2}\right)^{2}\left|\int_{0}^{L}\mathbf B_{\perp}(s)\,e^{iqs}\,ds\right|^{2}\equiv\left(\frac g2\right)^{2}K,
\qquad q=\frac{m_a^{2}-\omega_{\rm pl}^{2}}{2E},
\label{eq:conversion}
\end{equation}
with $\mathbf B_\perp=\mathbf B-(\mathbf B\cdot\hat n)\hat n$ a vector. The vector matters: where the transverse field rotates along the ray the components partly cancel, and an integral of $|\mathbf B_\perp|$ overstates the amplitude. The sign of the phase and the choice of transverse basis do not affect $P$. In SI units, $P=2.45\times10^{-21}\,(g/10^{-10}\ \mathrm{GeV^{-1}})^{2}\,(\sqrt K/\mathrm{T\,m})^{2}$, which reproduces the coefficient of \citet{yamamoto2020alp}.

For \DarkNESS{} every neglected term has a number. Weak mixing is exact at $P\le10^{-16}$. The ionospheric plasma frequency is below $3.7\times10^{-8}$~eV, so the plasma phase over $10\,R_E$ at 1~keV is below $2\times10^{-4}$ and $q=m_a^2/2E$. QED birefringence at $30\ \mu$T is $8\times10^{-33}$ of the phase and Faraday rotation is negligible at keV energies. The field is IGRF-14 to degree 13, internal only; the external magnetosphere is a separate model choice. A ray is integrated outward from the spacecraft along the line of sight to $10\,R_E$, where the dipole is $10^{-3}$ of its surface value, or to its first entry into the shell 150~km above the surface, below which the air is opaque to keV X-rays \citep{davoudiasl2006geomagnetic}. The photon keeps the ALP's direction, so only field in front of the detector along directions inside the aperture contributes; this is the fact on which the solar-axion interpretation of an XMM--Newton modulation failed \citep{fraser2014xmm,roncadelli2015nosolar}, and it works in favour of a source that fills the sky. The kernel is averaged over the $20^\circ$ cone with a 19-ray equal-area quadrature.

The implementation was checked against the uniform-field result $K=(\BpL)^2\,2(1-\cos qL)/(qL)^2$ to $10^{-5}$, against exact cancellation for a field that reverses at mid-path, against invariance under a joint rotation of field and ray, against convergence in step count on a limb-grazing ray, and against the Suzaku range $K=10^{4}$--$10^{5}\ \mathrm{T^2\,m^2}$ in Suzaku-like geometry.

\subsection{The source}

Let $m_\chi$ be the parent mass, $E_0=m_\chi/2$, $\tau$ the lifetime, and $f$ the parent's present-day share of the dark matter with the branching ratio to $aa$ folded in. Counting decays gives both components with one $f$ and one $\tau$ \citep{dror2021cosmic}. Parents along a line of sight with halo column $D(\hat n)=\int\rho_{\rm DM}\,ds$ decay now, so the Milky Way component is a line at $E_0$ of intensity
\begin{equation}
I_{\rm line}(\hat n)=\frac{f\,D(\hat n)}{2\pi\,m_\chi\,\tau},
\label{eq:line}
\end{equation}
Doppler-broadened by the halo dispersion to $\sigma_E/E_0=5.5\times10^{-4}$ \citep{evans2019shm}, 2~eV at 7~keV. Decays at scale factor $a$ emit daughters that arrive at $E=aE'$ and are diluted by $a^3$; with the Jacobian $dE'/dE=1/a$ and the survival of parents written in terms of today's share, the extragalactic component is
\begin{equation}
\frac{dI}{dE}=\frac{c}{4\pi}\,\frac{2f\,\rho_{{\rm DM},0}}{m_\chi\,\tau}\,\frac{e^{[t_U-t(z)]/\tau}}{E\,H(z)}\,\Theta(E_0-E),\qquad 1+z=\frac{E_0}{E},
\label{eq:continuum}
\end{equation}
with $H(z)$, $t(z)$ and $\rho_{{\rm DM},0}=1.257\ \mathrm{keV\,cm^{-3}}$ from Planck 2018 \citep{planck2018params}. Today's share is the convention the halo measures; a primordial share multiplies the line by $e^{-t_U/\tau}$, and the two agree only for $\tau\gg t_U$. For $\tau\gg t_U$ Eq.~\eqref{eq:continuum} reduces to the $E^{1/2}f(m_\chi/2E)$ power law of \citet{yamamoto2020alp} with $f(x)=(\Omega_m+\Omega_k/x+\Omega_\Lambda/x^3)^{-1/2}$; the sign of the last term is positive by the Friedmann equation, and with the sign as printed there $f(1)$ is imaginary.

The halo is NFW with $r_s=20$~kpc, $\rho_\odot=0.4\ \mathrm{GeV\,cm^{-3}}$ and $R_\odot=8.1$~kpc \citep{navarro1996nfw}. Its full-sky column integral, $2.67\times10^{23}\ \mathrm{GeV\,cm^{-2}\,sr}$, matches the value quoted by \citet{dror2021cosmic} to 1\%; the Milky Way and extragalactic energy densities stand in ratio 2.3 at long lifetime and their sum exceeds the CMB for $\tau$ below $10^{4}\,t_U$, as that work states.

\section{Counts in a skipper-CCD}
\label{sec:detector}

For time sample $i$ of length $\Delta t_i$ with cone-averaged kernel $K_i$, and energy bin $j$,
\begin{equation}
\mu^{\rm sig}_{ij}=\Delta t_i\left(\frac g2\right)^{2}K_i\,A_{\rm geo}\,\Omega\,f_{\rm live}\,\epsilon_{\rm sel}\int dE\,\frac{dI}{dE}(E)\,{\rm QE}(E)\,R_j(E),
\label{eq:counts}
\end{equation}
with each factor separate and applied once. The geometric area is 12~cm$^2$ and the $20^\circ$ full cone subtends $\Omega=0.0955$~sr, a grasp of 1.15~cm$^2$\,sr; the angular response is taken flat to the edge, an upper bound. The quantum efficiency is the transmission of the 50~nm aluminium window times the absorption in 500~$\mu$m of fully depleted silicon, from the NIST mass attenuation coefficients with their K edges: 0.984 at 1~keV, 0.951 at the Al edge, 0.981 at 10~keV. The redistribution $R_j$ is Gaussian with a width of 169~eV FWHM at 5.9~keV, the measured value \citep{alpine2026xray}, modelled as the Fano width ($F=0.118$, $w=3.72$~eV) and a 119~eV readout term in quadrature; the end-of-life width after three years in a Sun-synchronous orbit is 198~eV. Masking of charged-particle tracks leaves $f_{\rm live}=0.5$ \citep{alpine2025darkness}; the event-grade retention $\epsilon_{\rm sel}$ is unity here and falls with dose \citep{cervantesvergara2025proton,alpine2026xray}. For the Milky Way line the cone average is of the product $\langle DK\rangle_i$, which equals $\langle D\rangle\langle K\rangle$ to 0.5\% on the reference schedule.

The model was checked by a constant-brightness sky, which returns $I\,A\Omega\,\Delta E\,\Delta t$ exactly when every efficiency is unity; by conservation of counts through $R_j$ to $10^{-9}$; and by a narrow line that keeps its total and spreads more at 198~eV than at 169~eV. Both $f_{\rm live}$ and $\epsilon_{\rm sel}$ vary with the particle rate, which follows the same geomagnetic position as $K$; in a fit they are nuisance terms with their own time dependence.

\section{Likelihood and identifiability}
\label{sec:likelihood}

Counts are Poisson with mean
\begin{equation}
\mu_{ij}(\Th,\boldsymbol\beta)=\Th\,k_{ij}+\sum_p\beta_p\,z_{pij},\qquad \Th=\gag^{2}\frac{f}{\tau},
\label{eq:model}
\end{equation}
where $k_{ij}$ is the joint line-plus-continuum template per unit $\Th$ and the $z_p$ are background templates. The model is linear in $\Th$ because $P\propto g^2$ and both source components are $\propto f/\tau$; for $\tau\sim t_U$ the template keeps its $\tau$ dependence and $\Th$ is a label. The data constrain $\Th$; a value of $\gag$ follows only for a declared $f/\tau$, and the converted brightness per unit kernel is the source-agnostic quantity beneath both.

With Fisher information $I_{ab}=\sum_{ij}t_{a,ij}t_{b,ij}/\mu_{ij}$ over all amplitudes, the variance of $\hat\Th$ with every background amplitude free is $[I^{-1}]_{\Th\Th}$, the share of information that survives profiling is $1/(I_{\Th\Th}[I^{-1}]_{\Th\Th})$, and the 90\% one-sided limit is $\Th_{\rm UL}=1.28\,\sigma_\Th$, exact at $10^{4}$ counts per cell. Injected signals on a kernel-like template with a constant and a confounder are recovered with bias below $0.15\,\sigma$ and 1-$\sigma$ coverage between 62\% and 74\% over 400 Poisson realisations.

The reference schedule is a 420~km, $51.6^\circ$ orbit pointing at the Galactic Centre in umbra at 10-minute cadence, 39 science samples and 6.5~h per day, scaled to 187 days. The backgrounds are the cosmic X-ray background with the spectrum of \citet{deluca2004cxb} (9.9 counts~s$^{-1}$), a particle background scaling as the inverse cutoff rigidity with a flat spectrum and a mean level set equal to the CXB, an Earth-limb term $e^{-\theta_{\rm limb}/10^\circ}$ of unknown level entering as a shape, and a placeholder Galactic ridge model constant in time while the boresight is fixed. Across the science samples the kernel correlates with the limb term at 0.84 and with the particle proxy at $-0.67$.

\begin{table}[t]
\centering
\caption{Information on $\Th$ surviving the backgrounds, and the 90\% statistical limit, for the reference schedule over 187 days with $m_\chi=7$~keV. CXB plus particle background; the placeholder ridge model raises each limit by a factor 4 and changes no fraction by more than 0.1.}
\label{tab:identifiability}
\begin{tabular}{lcccc}
\toprule
Template & \multicolumn{2}{c}{Time--energy fit} & \multicolumn{2}{c}{Time-only fit}\\
 & kept & $\Th_{\rm UL}$ [GeV$^{-2}$\,Gyr$^{-1}$] & kept & $\sigma$ penalty\\
\midrule
joint & 0.945 & $3.8\times10^{-18}$ & 0.061 & $\times15$\\
line only & 0.954 & $3.9\times10^{-18}$ & 0.061 & $\times16$\\
continuum only & 0.559 & $1.4\times10^{-16}$ & 0.061 & $\times4.4$\\
\bottomrule
\end{tabular}
\end{table}

Table~\ref{tab:identifiability} gives the result. A fit to the time series alone keeps 6\% of the information on $\Th$, because the limb and cutoff templates track the kernel; a fit to the time--energy table keeps 95\%. The joint and line-only limits coincide and the continuum alone is 37 times weaker. With the spectrum in the fit the four nuisance templates remove 5\% of the information, so the central systematic of the method is the shape of the particle-background template, not its correlation with the kernel.

\section{Result}
\label{sec:results}

\subsection{The allowed region is out of reach}

The stellar bound $\gag<0.47\times10^{-10}\ \mathrm{GeV^{-1}}$ \citep{dolan2022globular} and the Planck 2018 bound on decaying dark matter, $f/\tau<4.0\times10^{-3}\ \mathrm{Gyr^{-1}}$ \citep{nygaard2021decaying,poulin2016decaying}, give $\Th_{\max}=8.8\times10^{-24}\ \mathrm{GeV^{-2}\,Gyr^{-1}}$. At $\Th_{\max}$ the reference day yields $2.4\times10^{-5}$ signal counts, of which $2.2\times10^{-5}$ are in the line, on $4.5\times10^{5}$ to $6.4\times10^{6}$ background counts; 187 days give $4.5\times10^{-3}$. The 90\% limit of Table~\ref{tab:identifiability} is $4\times10^{5}$ to $2\times10^{6}$ above $\Th_{\max}$; at the allowed $f/\tau$ it corresponds to $\gag=3$--$6\times10^{-8}\ \mathrm{GeV^{-1}}$ against the bound of $0.47\times10^{-10}$. A limit set by counting improves as the square root of exposure and grasp, so the gap corresponds to a factor $10^{11}$ in either; no schedule, aperture or lifetime choice provides it. The result is statistical, on idealised templates, and therefore a floor: the realistic limit is set by template error, as the Suzaku limit was.

\subsection{Against the Suzaku analysis}

The same machinery in Suzaku-like geometry (570~km, $31^\circ$, random blank-sky pointings, 12~Ms, grasp 0.026~cm$^2$\,sr, kernel median $1.4\times10^{4}\ \mathrm{T^2\,m^2}$) gives a 90\% statistical floor of $\gag=5.2\times10^{-8}\ \mathrm{GeV^{-1}}$ at $\tau=t_U$, or $9\times10^{-8}$ in the $3\sigma$ convention and the smaller grasp of \citet{yamamoto2020alp}. Their Eq.~4.1 at 7~keV is $2.1\times10^{-7}$, a factor 2.3 above the floor. In converted brightness per unit kernel their tabulated 2--6~keV limit sits 85 times above the floor; how that table is normalised per T$^2$\,m$^2$ is not resolved here and is carried as a caveat.

\subsection{What the analysis establishes}

Three statements carry the paper's contribution, each with a number and a test behind it. Identifiability is spectral: the energy spectrum, not the orbital modulation, separates the signal from what moves with it, and resolution with a spectral template is worth more than kernel contrast. The halo line carries the search, which favours a detector with the resolution to see a narrow feature at a known energy inside a wide field and a pointing that maximises the kernel-weighted halo column. The confounders cost little once modelled with the spectrum, so the systematic budget of an Earth-field search belongs to the particle-background template shape.

\section{Discussion}
\label{sec:discussion}

The placeholder ridge model and the assumed particle level move the limit by a factor 4 and no conclusion. The 150~km end of an occulted ray is a constant where the true opaque altitude rises with nadir angle and falls with energy; it affects only the frames used as a kernel-off control. The live fraction and selection efficiency, constants here, become time-dependent nuisance terms in a real fit. The front-side dead layer of the sensors and the angular response of the aperture are the largest declared unknowns and both act at 1--2~keV, where 40\% of the continuum counts fall. None of these reaches the $10^{5}$ in $\Th$ the gate would need; only a source outside the $\chi\to aa$ hypothesis, or a revision of the coupling or decay bounds, does.

Magnetic-field uncertainty is part of the signal model rather than a background \citep{matthews2022magneticmodels}; here the internal IGRF field is used alone, and the external magnetosphere is a separate model choice that the method can carry as a nuisance realisation. Lobster-eye optics remain a secondary trade: focusing does not by itself help a diffuse signal, and the result above says that what would help is resolution at the line, not grasp.

\section{Conclusions}
\label{sec:conclusions}

A \DarkNESS{}-class NanoSat cannot reach any $\chi\to aa$ ALP model allowed by current bounds through geomagnetic conversion, by five to six orders of magnitude in $\Th$, and no observing choice changes that. What it can contribute is the anatomy of the measurement: spectral identification, the halo line as carrier, and confounders that cost 5\% of the information once the spectrum is used. Those three statements are the design rules for any future Earth-field ALP search.

\section*{Data and code availability}

The analysis library, the state tables, the derivation notes with their checks, and the tests that pin every number in this paper are in the project repository. Archival identifiers will be added before submission.

\bibliographystyle{unsrtnat}
\bibliography{references,alp_dark_matter_30_new_references}

\end{document}
